\section*{NeurIPS Paper Checklist}

\begin{enumerate}

\item {\bf Claims}
    \item[] Question: Do the main claims made in the abstract and introduction accurately reflect the paper's contributions and scope?
    \item[] Answer: \answerYes{}
    \item[] Justification: The abstract and introduction state the controlled contamination study, lifecycle analyses, and evaluation-library correction. These claims are developed in Secs.~\ref{sec:methods}--\ref{sec:inference} and Appendix~\ref{app:sec:math_verify_bug}, with scope limitations discussed in Appendix~\ref{app:sec:discussion}.

\item {\bf Limitations}
    \item[] Question: Does the paper discuss the limitations of the work performed by the authors?
    \item[] Answer: \answerYes{}
    \item[] Justification: Appendix~\ref{app:sec:discussion} includes a limitations paragraph covering the single-benchmark focus, model scale and architecture, and corpus mixture. The main Discussion also points readers to that appendix.

\item {\bf Theory assumptions and proofs}
    \item[] Question: For each theoretical result, does the paper provide the full set of assumptions and a complete (and correct) proof?
    \item[] Answer: \answerNA{}
    \item[] Justification: The paper does not present formal theorems, lemmas, or proof claims. Its equations in Secs.~\ref{sec:pretraining} and~\ref{sec:inference} are empirical scaling-law fits and approximations used to interpret experiments.

    \item {\bf Experimental result reproducibility}
    \item[] Question: Does the paper fully disclose all the information needed to reproduce the main experimental results of the paper to the extent that it affects the main claims and/or conclusions of the paper (regardless of whether the code and data are provided or not)?
    \item[] Answer: \answerYes{}
    \item[] Sec.~\ref{sec:methods} and Appendix~\ref{app:sec:pretraining_implementation_details} describe the core recipe needed to reproduce the study, including data sources, contamination construction, model configurations, token budgets, optimizer settings, batch and learning-rate rules, sequence length, evaluation metrics, and training infrastructure.


\item {\bf Open access to data and code}
    \item[] Question: Does the paper provide open access to the data and code, with sufficient instructions to faithfully reproduce the main experimental results, as described in supplemental material?
    \item[] Answer: \answerNo{}
    \item[] Justification: Speculative draft, pending release audit: Appendix~\ref{app:sec:pretraining_implementation_details} links a public GitHub repository, names the MATH and FineWeb-Edu-Dedup data sources, and states that trained models were uploaded to HuggingFace Hub, so an open verification path may be intended. Finalize this as \answerYes{} only if the paper or supplement provides usable release links, environment and run instructions, data-preparation steps, and license or terms information for the released code, data products, and checkpoints; otherwise answer \answerNo{} and clearly identify what is not available.
    \item[] Guidelines:
    \begin{itemize}
        \item The answer \answerNA{} means that paper does not include experiments requiring code.
        \item Please see the NeurIPS code and data submission guidelines (\url{https://neurips.cc/public/guides/CodeSubmissionPolicy}) for more details.
        \item While we encourage the release of code and data, we understand that this might not be possible, so \answerNo{} is an acceptable answer. Papers cannot be rejected simply for not including code, unless this is central to the contribution (e.g., for a new open-source benchmark).
        \item The instructions should contain the exact command and environment needed to run to reproduce the results. See the NeurIPS code and data submission guidelines (\url{https://neurips.cc/public/guides/CodeSubmissionPolicy}) for more details.
        \item The authors should provide instructions on data access and preparation, including how to access the raw data, preprocessed data, intermediate data, and generated data, etc.
        \item The authors should provide scripts to reproduce all experimental results for the new proposed method and baselines. If only a subset of experiments are reproducible, they should state which ones are omitted from the script and why.
        \item At submission time, to preserve anonymity, the authors should release anonymized versions (if applicable).
        \item Providing as much information as possible in supplemental material (appended to the paper) is recommended, but including URLs to data and code is permitted.
    \end{itemize}


\item {\bf Experimental setting/details}
    \item[] Question: Does the paper specify all the training and test details (e.g., data splits, hyperparameters, how they were chosen, type of optimizer) necessary to understand the results?
    \item[] Answer: \answerYes{}
    \item[] Justification: Sec.~\ref{sec:methods} and Appendix~\ref{app:sec:pretraining_implementation_details} specify the model sizes, datasets, contamination levels, token budgets, optimizer, learning-rate and batch-size rules, sequence length, data mixing, and evaluation metrics. Secs.~\ref{sec:further_training} and~\ref{sec:inference} give the overtraining multipliers, SFT setup, sampling temperatures, and solution-length analyses.

\item {\bf Experiment statistical significance}
    \item[] Question: Does the paper report error bars suitably and correctly defined or other appropriate information about the statistical significance of the experiments?
    \item[] Answer: \answerNo{}
    \item[] Justification: Sec.~\ref{sec:methods} states that the study trains one model per parameter/replica configuration. The main figures report deterministic sweeps over model size, contamination level, overtraining, finetuning, temperature, and sequence length, but the paper does not report repeated-seed error bars, confidence intervals, or formal significance tests.

\item {\bf Experiments compute resources}
    \item[] Question: For each experiment, does the paper provide sufficient information on the computer resources (type of compute workers, memory, time of execution) needed to reproduce the experiments?
    \item[] Answer: \answerNo{}
    \item[] Justification: Secs.~\ref{sec:methods} and~\ref{sec:further_training} report token budgets and FLOP approximations, and Appendix~\ref{app:sec:pretraining_implementation_details} states that training used PyTorch DDP with NCCL. The paper does not report the hardware type, memory, wall-clock runtime, or total project compute needed to reproduce each experiment.
    
\item {\bf Code of ethics}
    \item[] Question: Does the research conducted in the paper conform, in every respect, with the NeurIPS Code of Ethics \url{https://neurips.cc/public/EthicsGuidelines}?
    \item[] Answer: \answerYes{}
    \item[] Justification: The authors have reviewed the NeurIPS Code of Ethics and confirm that the research conforms to it. The work consists of controlled model-training experiments on existing benchmarks and web-corpus data, with no crowdsourcing, human-subject experimentation, or deployment on users; anonymity is preserved in the submission, and any code, data, and model releases follow the data use and asset-release expectations of the Code of Ethics.


\item {\bf Broader impacts}
    \item[] Question: Does the paper discuss both potential positive societal impacts and negative societal impacts of the work performed?
    \item[] Answer: \answerYes{}
    \item[] Justification: Appendix~\ref{app:sec:broader_impacts} discusses both positive impacts (more trustworthy evaluations, simple model-agnostic checks for memorization, and a direct fix to a community evaluation harness) and negative impacts (inflated benchmark scores leading to brittle deployments, and the limited dual-use risk of describing memorization mechanisms). It also outlines mitigations practitioners can apply when probing models for contamination.
    
\item {\bf Safeguards}
    \item[] Question: Does the paper describe safeguards that have been put in place for responsible release of data or models that have a high risk for misuse (e.g., pre-trained language models, image generators, or scraped datasets)?
    \item[] Answer: \answerNA{}
    \item[] Justification: The study trains small research models, up to 344M parameters, to analyze benchmark contamination rather than to introduce a new capability model. Appendix~\ref{app:sec:broader_impacts} states that the work does not release new model capabilities, scraped data, or other assets that pose misuse risks independent of the contamination findings. The paper therefore does not describe a high-risk data or model release requiring controlled-access safeguards.
    \item[] Guidelines:
    \begin{itemize}
        \item The answer \answerNA{} means that the paper poses no such risks.
        \item Released models that have a high risk for misuse or dual-use should be released with necessary safeguards to allow for controlled use of the model, for example by requiring that users adhere to usage guidelines or restrictions to access the model or implementing safety filters. 
        \item Datasets that have been scraped from the Internet could pose safety risks. The authors should describe how they avoided releasing unsafe images.
        \item We recognize that providing effective safeguards is challenging, and many papers do not require this, but we encourage authors to take this into account and make a best faith effort.
    \end{itemize}

\item {\bf Licenses for existing assets}
    \item[] Question: Are the creators or original owners of assets (e.g., code, data, models), used in the paper, properly credited and are the license and terms of use explicitly mentioned and properly respected?
    \item[] Answer: \answerYes{}
    \item[] Justification: The paper cites or names the existing assets used in the study, including MATH, FineWeb-Edu-Dedup, Qwen 3, EleutherAI's LM Evaluation Harness, Flash Attention 2, PyTorch/DDP, and HuggingFace Hub. Appendix~\ref{app:sec:pretraining_implementation_details} explicitly lists the relevant public licenses or service terms for these assets.

\item {\bf New assets}
    \item[] Question: Are new assets introduced in the paper well documented and is the documentation provided alongside the assets?
    \item[] Answer: \answerTODO{}
    \item[] Justification: Speculative draft, pending asset-release confirmation: Appendix~\ref{app:sec:pretraining_implementation_details} says trained models were uploaded to HuggingFace Hub and links a public GitHub repository, so the submission may introduce code and checkpoint assets. Authors should change this to \answerYes{} only if the released assets include direct or anonymized links, licenses, intended-use and limitations notes, and model-card or repository documentation; if no new assets are released with the submission, \answerNA{} is the safer final answer.
    \item[] Guidelines:
    \begin{itemize}
        \item The answer \answerNA{} means that the paper does not release new assets.
        \item Researchers should communicate the details of the dataset\slash code\slash model as part of their submissions via structured templates. This includes details about training, license, limitations, etc. 
        \item The paper should discuss whether and how consent was obtained from people whose asset is used.
        \item At submission time, remember to anonymize your assets (if applicable). You can either create an anonymized URL or include an anonymized zip file.
    \end{itemize}

\item {\bf Crowdsourcing and research with human subjects}
    \item[] Question: For crowdsourcing experiments and research with human subjects, does the paper include the full text of instructions given to participants and screenshots, if applicable, as well as details about compensation (if any)? 
    \item[] Answer: \answerNA{}
    \item[] Justification: The experiments described in Secs.~\ref{sec:methods}--\ref{sec:inference} use existing benchmarks/corpora and model training/evaluation. The paper source does not describe crowdsourcing, participant studies, or human-subject data collection.

\item {\bf Institutional review board (IRB) approvals or equivalent for research with human subjects}
    \item[] Question: Does the paper describe potential risks incurred by study participants, whether such risks were disclosed to the subjects, and whether Institutional Review Board (IRB) approvals (or an equivalent approval/review based on the requirements of your country or institution) were obtained?
    \item[] Answer: \answerNA{}
    \item[] Justification: No crowdsourcing or human-subject experiments are described in the paper source. Therefore participant-risk disclosures and IRB or equivalent approvals do not appear applicable.

\item {\bf Declaration of LLM usage}
    \item[] Question: Does the paper describe the usage of LLMs if it is an important, original, or non-standard component of the core methods in this research? Note that if the LLM is used only for writing, editing, or formatting purposes and does \emph{not} impact the core methodology, scientific rigor, or originality of the research, declaration is not required.
    %this research? 
    \item[] Answer: \answerYes{}
    \item[] Justification: The paper's core experiments train and evaluate causal language models, and this usage is described in Sec.~\ref{sec:methods} and Appendix~\ref{app:sec:pretraining_implementation_details}. The source states that the models use the Qwen 3 architecture from random initialization and gives sizes, training setup, and evaluation metrics.

\end{enumerate}
